\section{Extensions to Deterministic Local Vote Rules}
\label{sec:extensions}

The main results of this paper are specific to deterministic $k$-NN.  However,
two aspects of the analysis generalize immediately to broader classes of
deterministic local vote rules.

\subsection{What generalizes}

\paragraph{Margin-crossing idea.}
Any deterministic local classifier whose prediction depends on a signed sum of
votes within a fixed-size neighborhood admits a margin-crossing criterion
analogous to \Cref{prop:5.1}.  If the prediction changes only when some
score crosses a threshold, then perturbation analysis reduces to tracking how
that score changes.  This covers weighted $k$-NN, distance-weighted local rules,
and certain decision-tree voting schemes.

\paragraph{Delete-plus-insert perturbation calculus.}
\Cref{prop:4.1} is already abstract: it applies to any learning rule with a
bounded loss.  The decomposition of replace-one into delete-one plus insert-one
is purely algebraic and does not depend on the classifier structure.  Any
stability analysis that establishes separate control over deletion and insertion
automatically obtains control over replacement.

\subsection{What remains $k$-NN-specific}

The explicit LOO/replace-one separation construction is $k$-NN-specific,
because it relies on ordered top-$k$ membership and duplicate occurrence
patterns that exploit distance ties at the query point.  The constructions in
\Cref{sec:odd-k-separation,sec:even-k-separation} specifically depend on the
ability to place multiple labeled copies at the same metric point and control
the order in which they enter the neighborhood.

For local rules without a hard top-$k$ cutoff (e.g., kernel rules with
decaying weights), the notion of replacing an occurrence may have a
quantitatively different effect, and the sharp separation exhibited here may
not directly transfer.

\subsection{Open characterization problem}

Which deterministic local vote rules admit an analogous LOO/replace-one
separation?  A full characterization would identify the necessary and
sufficient conditions on the vote aggregation function, the neighborhood
selection mechanism, and the tie-breaking protocol.

This is not a question the present paper answers.  Our contribution is to
establish that the separation exists for the most widely used local rule
($k$-NN), to provide a precise mechanism (margin crossing), and to record the
structural features that make it possible (ordered top-$k$ membership,
duplicate occurrences, deterministic tie-breaking).  Characterizing the
boundary of this phenomenon across the space of local learning rules remains an
open problem.
